\chapter*{Abstrak}
\addcontentsline{toc}{chapter}{Abstrak}

Rangkaian penderia Internet Pelbagai Benda (IoT) yang diagihkan sering beroperasi di bawah kekangan
tenaga yang ketat, terutamanya dalam penempatan bangunan pintar berskala besar. Kajian ini
mencadangkan rangka kerja penjadualan adaptif berasaskan pembelajaran mesin yang menyesuaikan
kadar persampelan penderia secara dinamik berdasarkan corak beban rangkaian dan variasi persekitaran.
Rangka kerja ini dinilai menggunakan set data dunia sebenar yang dikumpul daripada penempatan
bangunan pintar Nexa Systems di Kuala Lumpur, merangkumi 48 penderia sepanjang tempoh 90 hari.
Keputusan eksperimen menunjukkan pengurangan penggunaan tenaga sebanyak 27.4\% berbanding
penjadualan kadar tetap konvensional, dengan peningkatan kependaman data purata yang boleh diabaikan
(kurang daripada 120 milisaat). Sumbangan utama kajian ini ialah algoritma penjadualan ringan yang
sesuai untuk penggunaan pada get pinggir (edge gateway) dengan sumber terhad, serta metodologi
penilaian yang boleh disesuaikan untuk topologi rangkaian penderia yang lain.

\medskip
\noindent\textbf{Kata kunci:} Internet Pelbagai Benda, pembelajaran mesin, kecekapan tenaga,
pengkomputeran pinggir, penjadualan adaptif

% ---------- Abstract ----------
